%! ~~~ Packages Setup ~~~ 
\documentclass[]{../../../../tex/classes/styledArticle}
\usepackage{../../../../tex/packages/hebrewSupport}
\usepackage{../../../../tex/packages/mathShortcuts}
\usepackage{../../../../tex/packages/theoremsSupport}
\usepackage{../../../../tex/packages/enfancysections}

%! ~~~ Document ~~~

\author{שחר פרץ}
\title{\textit{ליניארית 2א 3}}
\date{31 למרץ 2025}
\begin{document}
	\maketitle
	\textbf{מרצה: }בן בסקין
	
	\defi{בעבור $A \in M_n(\F)$ הפולינום האופייני של $A$ הוא $f_A(x) = \det(Ix - A)$. }
	
	ראינו ש־$v$ ו"ע של $A$ עם ערך עצמי $\lg$ אמ"מ $v \in \ker(\lg I - A)$, וכן $\lg$ ע"ע אמ"מ $\dim \ker \lg - A > 0$. 
	
	\theo{$f_A(x)$ פולינםו מתוקן (מקדם מוביל 1) מדרגה $n$, המקדם של $x^{n - 1}$ הוא $-\tr A$, המקדם החופשי הוא $(-1)^{n}\det A$}. 
	
	\begin{proof}
		\begin{itemize}
			\item\textbf{ תקינות הפולינום. }מבין $n!$ המחוברים, ישנו אחד יחיד שדרגתו היא $n$. הסיבה היא שמדטרמיננטה לפי תמורות, התמורה היחידה שתתאים היא הזהות שתעבור על האלכסון. באינדוקציה על $n$, ונקבל: 
			\[ f_A(x) = (x - a_{11}) \cdots (x - a_{nn}) = x^{n} + \cdots \]
			באופן דומה אפשר להוכיח באינדוקציה באמצעות פיתוח לפי שורות: 
			\[ f_A(x) = (x - a_{11})|A_11| + \underbrace{a_{21}|A_21| - a_{31}|A_31| + \cdots + (-1)^{n + 1}a_{n1}|A_{n1}|}_{\text{מה.א. דרגה קטנה מ־$n$}} \]
			סה"כ גם כאן הראינו שהדרגה מתקבלת מהפולינום $\prod^{n}_{i = 1}(x - a_{ii})$
			
			\item\textbf{המקדם של $-\tr A$. }ההוכחה המסודרת מופיע העמוד 8 של הסיכום. מקדמי $x^{n - 1}$ מגיעים גם הם רק מ־* (הפולינום למעלה) שהם $-\tr A = \sum_{i = 1}^{n}-a_{ii}$ המקדם החופשי. 
			\item\textbf{המקדם החופשי. }$f_A(0) = \det(I \cdot 0 - A) = \det(-A) = (-1)^{n}\det A$. 
		\end{itemize}
	\end{proof}
	
	\subsection{הפסקה לרגע מליניארית 2א כדי לראות פתרון נורמלי רעיונית אך קצר לשאלה 3 ממועד ב' לליניארית 1א}
	יהי $V$ מ"ו $n$ ממדי, ו־$0 \neq v \in V$, וסקלרים $(\ag_i)_{i = 1}^{n}$. הוכיחו שקיים בסיס $(w_1 \dots w_n)$ כך ש־$v = \sum_{i = 1}^{n}\ag_iw_i$. 
	\begin{proof}הפתרון של גינסבורג: 
		אם בה"כ $\ag_1 \neq 0$ אז $w_1 := \ag_1\op\cl{v - \sum_{i = 2}^{n}\ag_iw_i}$ וכזה $(v, w_2 \dots w_n)$ בסיס וכאילו די גמרנו. 
	\end{proof}
	\begin{proof}פתרון נורמלי של המרצה: 
		נרצה למצוא $B_V = (v, v_w \dots v_n)$ כך שבהינתן $B_W = (w_1 \dots w_n)$ כך ש־: 
		\[ [v]_{B_V} = (1, 0, \dots 0), \ [v]_{B_{W}} = (\ag_1 \dots \ag_n) \]
		במילים אחרות, נגדיר איזו' מתאימה לשינוי בסיס. כלומר איזו $\phi \co V \to \F^n$ כך ש־$\phi(v) = (\ag_1, \dots, \ag_n)$ וכן $\phi(v_2) = e_2 \dots \phi(v_n) = e_n$ כך שסה"כ תיארנו העתקה שמזיזה בסיס אחד לאחר וסה"כ סיימנו. 
	\end{proof}
	
	\subsection{חזרה לקורס}
	\textbf{דוגמאות. }\begin{enumerate}[A)]
		\item אם $A = \binom{a \, b}{c \, d}$ אז $f_A(x) = x^{2} - (a + d)x + ad - bc$ (נטו מהמשפט הקודם). 
		\item אם $A = \diag(\lg_1 \dots \lg_n)$ אז $f_A(x) = \prod_{i = 1}^{n}(x - \lg_i)$. 
		\item אם $A = \pms{\lg_1 & & * \\ & \ddots & \\  0 && \lg_n}$ אז גם כאן $f_A(x) = \prod_{i = 1}^{n}(x - \lg_i)$ אך כדאי לשים לב שמשולשית עליונה לא בהכרח דומה לאלכסונית עם אותם הקבועים. 
		\item אם $A = \pms{B & * \\ 0 & C}$ כאשר $B, C$ בלוקים ריבועיים אז $f_A(x) = f_B(x) \cdot f_C(x)$. 
	\end{enumerate}
	
	\defi{בהינתן $T \co V \to V$} ט"ל נגדיר את הפולינום האופייני שלה (פ"א) באופן הבא: נבחר בסיס $B$ למ"ו $V$, ונתבונן ב־$A = [T]_B$ ונגדיר את $f_T(x) := f_A(x)$. 
	
	"אתה פותר עכשיו שאלה משיעורי הבית" "אל תדאג הבודק כבר שלח פתרון" "מה?!"
	
	\textbf{טענה. }הפ"א של ט"ל מוגדר היטב. למטריצות דומות אותו פ"א. \textit{הוכחה. }בשיעורי הבית. ויש סיכוי שגם בדף הסיכום. 
	
	\textit{מטלה. }לקרוא בעמוד 10 על שוויון פונ' בין שדות שונים, רלוונטי כדי להראות שלא משנה על איזה שדה אנו מדברים הפולינום האופייני נשאר זהה. 
	
	\textbf{דוגמה. }נתבונן בהעתקה $\R_n[x] \to \R_n[x], \ T(f) = f'$. נבחר בסיס $B = (1, x, \dots, x^{n})$. אז: 
	\[ [T]_B = \pms{0 & 1  & 0 & 0 & \dots & 0 \\ 0 & 0 & 2 & 0 & \dots & 0 \\ 0 & 0 & 0 & 3 & \ddots & \vdots \\ \vdots & \vdots & \vdots & \ddots & \ddots \\ 0 & 0 & 0 & 0 } \]
	אז: 
	\[ f_T(x) = \det\pms{x & -1 & 0 &\dots \\ & x & -2 & 0 &\dots \\ && \ddots &\ddots } = x^{n + 1} \]
	\theo{$T \co V \to V$ ט"ל, אז $\lg$ ע"ע של $T$ אמ"מ $f_T(\lg) = 0$. }
	\begin{proof}
		יהא $B \subseteq V$ בסיס של $V$. אז $A = [T]_B$ ואז $f_T(\lg) = 0$ אמ"מ $f_A(\lg) = 0$ אמ"מ $\lg$ ע"ע של $A$. 
	\end{proof}
	
	\textbf{הגדרה. }יהי $\lg \in \F$ ע"ע של $T$ (או $A$). \textit{האיברוי האלגברי} של $\lg$ הוא החזקה המקסימלית $d$ כך ש־$(x - \lg)^{d} \mid f_T(x)$ (חלוקת פולינומים). 
	
	\textbf{דוגמה. }בעבור $T$ היא הגזירה, ממקודם: $f_T(x) = x^{n + 1}$ ע"ע יחיד הוא $0$. הריבוי האלגברי של $0$ הוא $n + 1$. הריבוי הגיאומטרי של $0$ הוא $1$. 
	
	\noti{נניח ש־$\lg$ ע"ע של $T$ (או $A$) אז $d_\lg$ הריבוי האלגברי של $\lg$ ו־$r_\lg$ הריבוי הגיאומטרי של $\lg$. }
	\theo{$r_\lg \le d_\lg$ ושוויון אמ"מ $T$ (או $A$) לכסינה. }
	
	\ndoc
\end{document}
